% Predictability probe: the in-band generator-inversion attack, measured.
% Standalone fragment; evidence in paper/evidence/predictability.json, figure in
% paper/figures/predictability.pdf, script paper/src/make-predictability.py.

\subsection{Predictability probe: measuring the in-band inversion attack}\label{sec:predictability}

\Cref{sec:limitations} names a threat the leak-free interface property does not cover: a scenario is a pure function of a public 64-bit seed under a public deterministic transform, so future bars are in principle computable from past bars, and no forbidden operation is needed. This probe converts that named threat into a measurement. Three adversaries predict the next bar's return from a causal prefix only, on the three volatility tiers, 16 seeds per tier (4 symbols, 120 days, 30 warmup bars): an \emph{unconditional baseline} that predicts the expanding prefix-mean return per symbol; an \emph{honest statistical adversary}, a ridge AR(5) refit on the prefix at every bar, with no knowledge of the generator; and a \emph{known-seed oracle}, which replays the generator from the true seed and therefore predicts every future bar exactly, the upper bound on any inversion attack. Each predictor is also run as a frictionless unit-gross long/short sign-following policy and scored through the public \texttt{score\_run} kernel; we report the kernel's deflated-Sharpe probability.

\paragraph{The honest adversary extracts almost nothing.} On Calm the ridge AR reaches $55.4\% \pm 2.8\%$ directional accuracy against the baseline's $54.1\% \pm 2.4\%$: it detects the tier's designed momentum autocorrelation, worth about one accuracy point over drift-following, and its policy DSR ($0.97 \pm 0.06$ vs.\ $0.90 \pm 0.12$) reflects the same mild, real, and disclosed structure rather than any transform inversion. On the stressed tiers the edge vanishes: $51.2\% \pm 3.3\%$ on Hard and $50.6\% \pm 2.1\%$ on Extreme, statistically indistinguishable from coin flipping, with MSE no better than the unconditional baseline (Hard: $1.24\times10^{-4}$ vs.\ $1.22\times10^{-4}$; Extreme: $1.27\times10^{-3}$ vs.\ $1.17\times10^{-3}$) and policy DSRs of $0.48 \pm 0.35$ and $0.35 \pm 0.24$. The amplified tiers behave as near-martingales to a regression adversary: in-band statistical prediction recovers the disclosed stylized structure and nothing beyond it.

\paragraph{The oracle bound is total, and a bounded seed search reaches it.} The known-seed oracle scores $100\%$ accuracy, zero MSE, and DSR $1.00$ on every tier: determinism means the gap between an honest adversary and a successful inverter is the entire distance from coin flipping to omniscience, so the question is only the cost of inversion. We implemented the strongest cheap version: a brute scan over a bounded seed band of width $2^{16}$ centered on the published \texttt{start\_level}, matching each candidate's first-bar closes against the single observed opening bar. The scan recovered the true seed in 16 of 16 trials with zero collisions, from a prefix of one bar, at $14.7\,\mu$s per candidate ($\approx 0.96$\,s per scan); every recovered seed was then verified to reproduce the full 30-bar observed prefix on its tier. The lesson is sharp: when the deployed seed band is bounded and public, as a \texttt{(start\_level, num\_levels)} spec makes it, in-band inversion is not cryptanalysis but a table scan, and any such deployment should treat its scenarios as fully known to a motivated adversary. Extrapolating the measured per-candidate cost to the unbounded $2^{64}$ space gives $\sim 8.6$ million CPU-years, so brute-force inversion of an unpublished seed drawn from the full space is infeasible by enumeration.

\begin{figure}[t]
\centering
\includegraphics[width=0.78\linewidth]{predictability}
\caption{Next-bar directional accuracy by tier. The causal AR adversary remains near the unconditional baseline; knowing or recovering the seed gives perfect prediction.}
\label{fig:predictability}
\end{figure}

\paragraph{What remains open.} The probe brackets the attack but does not close it. Between the honest regression (no edge beyond disclosed structure) and the band scan (total recovery given a bounded public band) lies the case that matters for unbounded deployments: analytic inversion of SplitMix64 through the mul/add/div/max price transform from observed reals. SplitMix64 is not a cryptographic generator; its state transition is a fixed odd-constant addition and its output mixing is invertible, so recovering the state from a handful of \emph{exact} raw outputs is straightforward. What the adversary observes here is not raw outputs but prices, each a truncated 53-bit draw passed through data-dependent arithmetic with floors, and whether that quantization and composition resists lattice or guess-and-determine state recovery is a cryptanalytic question this paper does not answer. The honest reading of the probe: the environment's unpredictability rests on the secrecy and entropy of the seed, not on any hardness of the transform, and published bounded seed bands forfeit both.
